Transition-metal clusters offer useful models for examining how composition,
geometry, spin state, and local coordination influence adsorption at catalytic
metal sites. Their finite size gives access to atom-specific environments that
are difficult to isolate on extended surfaces, while still retaining key
features of heterogeneous catalysts such as low coordination, structural
fluxionality, and strong coupling between geometry and electronic structure.
Platinum-based materials are particularly relevant in this context because of
their high activity in fuel-cell-related processes and in a wide range of
catalytic reactions.\cite{jimenezizal:2018,jimenezizal:2019b,yu:2002,
zhang:2014,auer:1998} This performance, however, comes with a well-known
limitation, since Pt binds carbon monoxide strongly. Even small amounts of \ce{CO}
can block Pt active sites and cause severe deactivation, especially in
fuel-cell environments.\cite{cheng:2007,vogel:1975,baschuk:2001,
christoffersen:2001,yang:2021} Alloying Pt with a second element is therefore
a common strategy to tune the electronic structure of the active sites, reduce
Pt loading, and mitigate \ce{CO} poisoning.\cite{yousaf:2016,yin:2011,
stolbov:2009,lee:1999,hayden:2003,samjeske:2002,ferrari:2016,ahmad:2015,
ehteshami:2013,chen:2019,ribeiro:2019}

Germanium has emerged as a particularly interesting promoter for Pt. Pt--Ge
materials have been reported to improve activity or durability in hydrogen
evolution, direct methanol fuel cells, \ce{CO} oxidation, dehydrogenation, and
reforming reactions, with the Ge component modifying both the structure and the
electronic properties of the Pt sites.\cite{lim:2018,veizaga:2014,
veizaga:2019,crabb:2001,gyorffy:2009,poths:2023,vilella:2008,rimaz:2019,
jimenezizal:2019,ma:2023,garetto:1996,mazzieri:2009,mariscal:2007}
Experimental and theoretical studies on \ce{PtGe/C} electrocatalysts have also
shown modified \ce{CO} oxidation behavior relative to monometallic Pt
systems.\cite{veizaga:2014,veizaga:2019,crabb:2001} These observations suggest
that Ge can act through both geometric effects, by disrupting contiguous Pt
ensembles, and electronic effects, by changing the electron density and orbital
availability at Pt.
Experiments on \ce{Pt(111)-Ge} surface alloys also reported weaker \ce{CO}
chemisorption than on \ce{Pt(111)} and associated this change with the modified
electronic structure of Pt after Ge incorporation.\cite{fukutani:1996,fukutani:1998}
Earlier calculations on Pt--Ge cluster models likewise found weaker Pt--C
interactions and changes in the \ce{C-O} stretching frequency after Ge
substitution.\cite{liao:2000}

Gas-phase and supported clusters provide a controlled route to examine these
effects in atomically defined environments. A recent review of \ce{CO}
adsorption on doped metal clusters emphasizes the coupled roles of composition,
charge state, and local environment.\cite{ferrari:2020} For 13-atom Pt--Au
clusters, Morrow and co-workers showed that the local environment of the
adsorption atom affects both the \ce{CO} binding energy and stretching
frequency.\cite{morrow:2011} Ugartemendia and co-workers combined density
functional theory (DFT)
with mass spectrometry to study \ce{Pt_n^+} and \ce{GePt_{n-1}^+} clusters
with $n=5$--9.\cite{ugartemendia:2021} Their global-minimum searches and
electronic-structure analyses showed that replacing one Pt atom by Ge generally
weakens the Pt--CO interaction while the clusters remain active toward
\ce{H2} dissociation. Within a Blyholder-type interpretation, the weaker
{\ce{CO} binding was assigned mainly to reduced Pt $\rightarrow$ \ce{CO}
$\pi$ back-donation. In this interpretation, Pt--Ge orbital mixing competes
with the Pt orbitals that would otherwise interact with the \ce{CO} $2\pi^*$
acceptor manifold.}

A subsequent study extended this idea from monodoped cationic clusters to
\ce{Pt_{n-m}Ge_m} nanoclusters with different sizes and
compositions.\cite{ugartemendia:2022} {That work found especially weak
\ce{CO} binding and high resistance to poisoning in clusters with roughly
equimolar Pt:Ge ratios. The authors traced this behavior to strong Pt--Ge
covalency and electrostatic stabilization from favorable Pt--Ge mixing.} In
the same work, AdNDP analysis revealed a
clear bonding contrast between pure and equimolar clusters. Pure \ce{Pt10}
was described by delocalized multicenter Pt bonding, whereas \ce{Pt5Ge5}
could be represented mainly by localized Pt--Ge two-center bonds. {This
contrast suggests that weaker \ce{CO} adsorption may depend on the Ge content
and on how Ge changes the bonding framework around the Pt adsorption site.}

More recent work on supported \ce{PtGe} clusters has reinforced this picture
while adding the role of the support and reaction environment. Deposited
\ce{Pt4Ge_n} clusters on MgO were proposed to operate as bifunctional
catalysts for \ce{CO} oxidation, with \ce{CO} binding preferentially at Pt and
oxygen binding at Ge-derived sites; at the same time, Ge alloying reduces the
affinity of Pt for \ce{CO} and relieves self-poisoning.\cite{ugartemendia:2024}
Related studies in alkane dehydrogenation further indicate that the Pt:Ge ratio
can control activity, selectivity, coke resistance, and sintering resistance in
supported Pt--Ge clusters.\cite{poths:2023b} Taken together, these studies show
that Ge is not a passive diluent. It changes the local chemistry of Pt, the
stability of the cluster, and the possible reaction pathways available under
catalytic conditions.

The present work examines ten-atom \ce{Pt10}, \ce{Pt_{x}Ge_{10-x}}, (where x=4,5,6),
clusters using DFT\@. The study uses two crossed
construction routes. In the first, we start from the previously reported
lowest-energy, or global-minimum (GM), \ce{Pt10} framework and replace selected
Pt atoms by Ge. In the reverse route, we start from the previously reported GM
\ce{Pt5Ge5} framework and replace all Ge atoms by Pt, producing a pure-
\ce{Pt10} derivative.\cite{ugartemendia:2022} The resulting structures provide
selected comparisons of composition and starting geometry within related
ten-atom frameworks.

\ce{CO} is then placed in turn at each independent Pt environment represented
in these parent structures, and each structure is fully optimized. This
procedure shows how the binding energy depends on composition, parent geometry,
spin state, local bonding, and relaxation of the whole cluster. In addition, we
start separate calculations from the reported global-minimum structure of each
composition. Comparing the two sets shows whether the adsorption pattern changes
when a different parent structure or spin state is used.
